%!TEX program = xelatex
\def\ano{2020}
\def\semestre{1}
\def\disciplina{\'Algebra 1}
\def\turma{C}
\def\autor{Jos\'e Ant\^onio O. Freitas}
\def\instituto{MAT-UnB}

\documentclass{beamer}
\usetheme{Madrid}
\usecolortheme{beaver}
% \mode<presentation>
\usepackage{caption}
\usepackage{textpos}
\usepackage{amssymb}
\usepackage{amsmath,amsfonts,amsthm,amstext}
\usepackage[brazil]{babel}
% \usepackage[latin1]{inputenc}
\usepackage{graphicx}
\graphicspath{{/home/jfreitas/GitHub_Repos/video_aulas/logo/}{C:/Users/josea/Documents/GitHub_Repos/video_aulas/logo/}}
\usepackage{enumitem}
\usepackage{multicol}
\usepackage{answers}
\usepackage{tikz,ifthen}
\usetikzlibrary{lindenmayersystems}
\usetikzlibrary[shadings]

\renewcommand{\qedsymbol}{$\blacksquare$}
\newtheorem{definicao}{Defini\c{c}\~ao}[section]
\newtheorem{definicoes}{Defini\c{c}\~oes}[section]
\newtheorem{exemplo}{Exemplo}[section]
\newtheorem{exemplos}{Exemplos}[section]
\newtheorem{exercicio}{Exerc{\'\i}cio}
\newtheorem{observacao}{Observa{\c c}{\~a}o:}[section]
\newtheorem{observacoes}{Observa{\c c}{\~o}es:}[section]
\newtheorem*{solucao}{Solu{\c c}{\~a}o:}
\newtheorem{proposicao}{Proposi\c{c}\~ao}
\newtheorem{corolario}{Corol\'ario}
\newenvironment{prova}[1][Prova]{\noindent\textbf{#1:} }{\qedsymbol}%{\ \rule{0.5em}{0.5em}}
\newcommand{\nsub}{\varsubsetneq}
\newcommand{\vaz}{\emptyset}
\newcommand{\im}{{\rm Im\,}}
\newcommand{\sub}{\subseteq}
\newcommand{\n}{\mathbb{N}}
\newcommand{\z}{\mathbb{Z}}
\newcommand{\rac}{\mathbb{Q}}
\newcommand{\real}{\mathbb{R}}
\newcommand{\complex}{\mathbb{C}}
\newcommand{\cp}[1]{\mathbb{#1}}
\newcommand{\ch}{\mbox{\textrm{car\,}}\nobreak}
\newcommand{\vesp}[1]{\vspace{ #1  cm}}
\newcommand{\compcent}[1]{\vcenter{\hbox{$#1\circ$}}}
\newcommand{\comp}{\mathbin{\mathchoice
{\compcent\scriptstyle}{\compcent\scriptstyle}
{\compcent\scriptscriptstyle}{\compcent\scriptscriptstyle}}}

\title{An\'eis - Suban\'eis}
\author[\autor]{\autor}
\institute[\instituto]{\instituto}
\date{}

\begin{document}
    \begin{frame}
        \maketitle
    \end{frame}

    \logo{\includegraphics[scale=.1]{logo-MAT.png}\vspace*{8.5cm}}


    \begin{frame}
        \vspace{1cm}
        $A \ne \emptyset$, $\oplus$ \pause e $\otimes$ operações binárias. \pause $(A, \oplus, \otimes)$ \pause {\'e} um \textbf{anel} se:\pause
        \begin{enumerate}[label={\roman*})]
            \item \pause para todos $x$, \pause $y$, \pause $z \in A$ \pause vale\pause
            \[
                (x \oplus y) \pause \oplus z \pause = x \oplus \pause (y \oplus z).\pause
            \]

            \item \pause Para todos $x$, \pause $y \in A$ \pause vale\pause
            \[
                x \oplus y = \pause y \oplus x.\pause
            \]

            \item \pause Existe $0_A \in A$ \pause tal que para todo elemento $x \in A$ \pause vale\pause
            \[
                x \oplus 0_A \pause = x \pause = 0_A \oplus x.\pause
            \]

            \item \pause Para cada elemento $x \in A$, \pause existe $y \in A$ \pause tal que\pause
            \[
                x \oplus y \pause = 0_A \pause = y \oplus x.\pause
            \]

            \seti
    \end{enumerate}
    \end{frame}

    \begin{frame}
        \begin{enumerate}[label={\roman*})]
            \conti

            \item \pause Para todos $x$, \pause $y$, \pause $z \in A$, \pause vale\pause
            \[
                (x\otimes y) \pause \otimes z \pause = x \otimes \pause (y\otimes z).\pause
            \]

            \item Para todos $x$, \pause $y$, \pause $z \in A$ \pause vale\pause
            \[
                (x \oplus y) \pause \otimes z \pause = x \otimes z \pause \oplus \pause y\otimes z.\pause
            \]

            \item Para todos $x$, \pause $y$, \pause $z \in A$ \pause vale\pause
            \[
                x \otimes \pause (y \oplus z) \pause = x \otimes y \pause \oplus \pause x\otimes z.\pause
            \]
        \end{enumerate}
    \end{frame}


    \begin{frame}
        \begin{observacao}
            Seja $(A, \oplus, \otimes)$ \pause um anel. \pause Para
            simplificar a nota\c{c}\~ao \pause vamos denotar a
            opera\c{c}\~ao $\oplus$ \pause por $+$ \pause e a
            opera\c{c}\~ao $\otimes$ \pause por $\cdot$ \pause e
            assim escrever simplesmente \pause que $(A, +, \cdot)$
            \pause \'e um anel.\pause
        \end{observacao}
    \end{frame}

    \begin{frame}
        \begin{proposicao}
            Seja $(A, + , \cdot)$ um anel. \pause Ent\~ao:\pause
            \begin{enumerate}[label={\roman*})]
                \item O elemento neutro {\'e} {\'u}nico.\pause

                \vspace{.5cm}

                \item Para cada $x \in A$ \pause existe um {\'u}nico
                    oposto. \pause Neste caso o \textbf{oposto} de
                    $x \in A$ \pause ser\'a denotado por $-x$.\pause

                \vspace{.5cm}

                \item Para todo $x \in A$, \pause
                \[
                    -(-x) = x.\pause
                \]

                \vspace{.5cm}

                \item Dados $x_{1}$, \pause $x_{2}$, \pause \dots,
                    $x_n \in A$, \pause $n \geqslant 2$, \pause
                    ent{\~a}o\pause
                \[
                    -(x_1 + x_2 + \dots + x_n) \pause = (-x_1) \pause
                    + (-x_2) \pause + \dots + (-x_n).\pause
                \]

                \vspace{.2cm}

                \seti
            \end{enumerate}
        \end{proposicao}
    \end{frame}

    \begin{frame}
        \begin{proposicao}
            \begin{enumerate}[label={\roman*})]
                \conti

                \item Para todos $\alpha$, \pause $x$, \pause $y \in A$, \pause se
                \[
                    \alpha + x \pause = \alpha + y,\pause
                \]
                ent{\~a}o $x = y$.\pause

                \vspace{.5cm}

                \item Para todo $x \in A$, \pause
                \[
                    x\cdot 0_A \pause = 0_A \pause = 0_A\cdot x.\pause
                \]

                \vspace{.5cm}
            \end{enumerate}
        \end{proposicao}
    \end{frame}

    \begin{frame}
        \begin{proposicao}
            \begin{enumerate}
                \item[vii)] Para todos $x$, \pause $y \in A$, \pause temos\pause
                \[
                    x \cdot (-y) \pause = (-x) \cdot y \pause = -(x \cdot y).\pause
                \]

                \vspace{.5cm}

                \item[viii)] Para todos $x$, \pause $y \in A$, \pause
                \[
                    x \cdot y \pause = (-x) \cdot (-y).\pause
                \]

                \vspace{.5cm}
            \end{enumerate}
        \end{proposicao}
    \end{frame}

    \begin{frame}
        \begin{definicao}
            Seja $(A, +, \cdot)$ um anel. \pause Dizemos que um subconjunto n{\~a}o vazio \pause $B\subseteq A$ \pause {\'e} um \textbf{subanel} de $A$ \pause quando $(B, +, \cdot)$ \'e um anel.\pause
        \end{definicao}

        \begin{exemplos}
            \begin{enumerate}[label={\arabic*})]
                \item Todo anel $A$ sempre tem dois suban{\'e}is: \pause $\{0_{A}\}$ \pause e $A$, \pause que s{\~a}o chamados de \textbf{suban{\'e}is triviais}.\pause

                \vspace{.5cm}

                \item Em $(\z_4,\oplus,\otimes)$ \pause o conjunto $B = \{\overline{0}, \overline{2}\}$ \'e um subanel.\pause

                \vspace{.5cm}

                \item No anel $\z$, \pause o conjunto $m\z = \{mk \mid k \in \z\}$, $m > 1$ \pause {\'e} um subanel de $\z$.\pause

                \vspace{.5cm}
            \end{enumerate}
        \end{exemplos}
    \end{frame}

    \begin{frame}
        \begin{proposicao}
            Seja $(A, +,\cdot)$ um anel. \pause Um subconjunto n{\~a}o vazio \pause $B\subseteq A$ \pause {\'e} um subanel de $A$ \pause se, e somente se, \pause $x - y = x + (-y) \in B$ \pause e $x\cdot y \in B$ \pause para todos $x$, $y \in B$.\pause
        \end{proposicao}
    \end{frame}

    \begin{frame}
        \begin{exemplos}
            \begin{enumerate}
                \item[1)] Em $(\z_4,\oplus,\otimes)$ \pause o conjunto $B = \{\overline{0}, \overline{2}\}$ \'e um subanel.\pause

                \vspace{5cm}
            \end{enumerate}
        \end{exemplos}
    \end{frame}

    \begin{frame}
        \begin{exemplos}
            \begin{enumerate}
                \item[2)] No anel $\z$, \pause o conjunto $m\z$, $m > 1$ \pause {\'e} um subanel de $\z$.\pause

                \vspace{5cm}
            \end{enumerate}
        \end{exemplos}
    \end{frame}

    \begin{frame}
        \begin{exemplos}
            \begin{enumerate}
                \item[3)] Considere o anel $(\rac, \star, \odot)$ \pause onde as opera\c{c}\~oes $\star$ e $\odot$ são definidas por\pause
                \begin{center}
                    $x \star y = x + y - 8$\\\pause
                    $x \odot y = x + y - \dfrac{xy}{8}.$\pause
                \end{center}
                Quais dos seguintes subconjuntos s\~ao suban\'eis?\pause
                \begin{enumerate}
                    \item[(a)] $B = \{2k \mid k \in \z\}$\pause

                    \vspace{.5cm}

                    \item[(b)] $C = \{8k \mid k \in \z\}$\pause
                    \vspace{.5cm}
                \end{enumerate}
            \end{enumerate}
        \end{exemplos}
    \end{frame}
\end{document}
